% BHU PYQ -- Manually transcribed & error-corrected
\documentclass[11pt,a4paper]{article}

\usepackage[a4paper,top=2.5cm,bottom=2.5cm,left=2.8cm,right=2.8cm,headheight=14pt]{geometry}
\usepackage{amsmath,amssymb}
\usepackage[T1]{fontenc}
\usepackage[utf8]{inputenc}
\usepackage{lmodern}
\usepackage{microtype}
\usepackage{fancyhdr}
\usepackage{enumitem}
\usepackage{hyperref}
\usepackage{lastpage}
\usepackage{parskip}

\hypersetup{colorlinks=true,urlcolor=black,linkcolor=black,
  pdftitle={CHB-605: Atomic & Molecular Structure -- BHU PYQ},pdfauthor={Banaras Hindu University}}

\pagestyle{fancy}\fancyhf{}
\renewcommand{\headrulewidth}{0.4pt}\renewcommand{\footrulewidth}{0.4pt}
\fancyhead[L]{\small\textit{Banaras Hindu University}}
\fancyhead[R]{\small\textit{CHB-605: Atomic \& Molecular Structure}}
\fancyfoot[C]{\small Page \thepage\ of \pageref{LastPage}}

\newlist{parts}{enumerate}{2}
\setlist[parts,1]{label=(\alph*),leftmargin=*,itemsep=5pt,topsep=3pt}
\setlist[parts,2]{label=(\roman*),leftmargin=*,itemsep=3pt,topsep=2pt}
\newcommand{\pts}[1]{\hfill\textit{[#1]}}

\begin{document}

\begin{center}
  {\large\bfseries BANARAS HINDU UNIVERSITY}\\[2pt]
  {\normalsize B.Sc. (Hons.) Semester VI Examination, 2022-23}\\[6pt]
  \rule{0.8\linewidth}{0.5pt}\\[6pt]
  {\large\bfseries Paper No. CHB-605: Atomic and Molecular Structure and Spectroscopic Techniques}\\[4pt]
  {\normalsize Time: 3 Hours\hfill Maximum Marks: 70}
\end{center}
\rule{\linewidth}{0.4pt}
\small
\noindent\textit{Instructions: Answer five questions in total, including Question~1 which is compulsory.}
\normalsize
\bigskip

\noindent\textbf{Question 1.} Answer all parts of the following: \pts{5 $\times$ 2 = 10}
\begin{parts}
  \item Deduce the point group and give one example of a molecule that is linear and centrosymmetric ($\text{D}_{\infty\text{h}}$).
  \item Write the Schr\"odinger wave equation for the Helium atom.
  \item State the Born-Oppenheimer approximation.
  \item State the Beer-Lambert law.
  \item Why is TMS (tetramethylsilane) used as an internal reference in NMR spectroscopy?
\end{parts}

\medskip\rule{\linewidth}{0.2pt}

\noindent\textbf{Question 2.}
\begin{parts}
  \item Draw the molecular orbital energy level diagram for the hydrogen molecule ($\text{H}_2$). Compare it with the Valence Bond wave function. \pts{6}
  \item Deduce the point group and give one example of a molecule that is non-linear and has $\text{C}_{2\text{v}}$ symmetry. \pts{8}
\end{parts}

\medskip\rule{\linewidth}{0.2pt}

\noindent\textbf{Question 3.}
\begin{parts}
  \item Explain why the $\lambda_{\max}$ of the $\text{n} \to \pi^*$ transition in an $\alpha,\beta$-unsaturated carbonyl compound shifts to shorter wavelengths (hypsochromic shift) with increasing solvent polarity. \pts{6}
  \item Sketch and interpret the $^1\text{H}$-NMR spectra of: (i) Ethyl bromide and (ii) Acetaldehyde, showing spin-spin splitting patterns. \pts{8}
\end{parts}

\medskip\rule{\linewidth}{0.2pt}

\noindent\textbf{Question 4.}
\begin{parts}
  \item Sketch the mass spectrum of $1$-bromopropane, showing the characteristic isotopic pattern. How can alkyl chlorides and bromides be distinguished by mass spectrometry? \pts{6}
  \item Predict the expected $\lambda_{\max}$ using Woodward-Fieser rules for the following compounds: \pts{8}
  \begin{parts}
    \item $4$-methyl-$3$-penten-$2$-one
    \item Cholesta-$3,5$-diene
  \end{parts}
\end{parts}

\medskip\rule{\linewidth}{0.2pt}

\noindent\textbf{Question 5.} On the basis of the following spectroscopic data, suggest a structure for the compound ($\text{C}_3\text{H}_7\text{NO}$): \pts{14}
\begin{itemize}
  \item $\text{IR} \ (v_{\max}): 3428\,\text{cm}^{-1}\text{ (m)}, 2941\text{--}2857\,\text{cm}^{-1}\text{ (w)}, 1681\,\text{cm}^{-1}\text{ (s)}, 1550\,\text{cm}^{-1}\text{ (m)}, 1452\,\text{cm}^{-1}\text{ (w)}$.
  \item $^1\text{H-NMR} \ (\delta, \text{ppm}): 2.02 \ (s, 3\text{H}), 2.70 \ (s, 3\text{H}), 8.13 \ (s, 1\text{H})$.
  \item $^{13}\text{C-NMR} \ (\delta, \text{ppm}): 17.7, 26.5, 171.2$.
\end{itemize}

\vfill
\rule{\linewidth}{0.4pt}
\begin{center}\small
  \href{https://bhu-exam.vercel.app/}{Click here to visit our website}\\[4pt]
  \href{https://me-aryan.vercel.app/}{Click here to visit our contributor}\\[4pt]
  \href{https://quashan-me.vercel.app/}{Click here to visit our contributor}
\end{center}

\end{document}
